\item \textbf{Diseño de cuencas de agua y resonancia de terremotos (seiches).}
Está trabajando como asistente de un arquitecto paisajista, que diseña los jardines alrededor de un nuevo edificio comercial. El arquitecto planea tener una gran cuenca de agua rectangular como parte de su diseño. Cuando ve este diseño, menciona al arquitecto que el proyecto está ubicado en un área propensa a los terremotos. Usted señala que un terremoto podría crear un seiche en la cuenca por resonancia, haciendo que el agua en la cuenca se derrame e ingrese a los transformadores eléctricos subterráneos cercanos.

Un seiche es una onda estacionaria en un cuerpo de agua, en el que el agua se balancea hacia adelante y hacia atrás con antinodos en los extremos de la cuenca. De acuerdo con datos empíricos tomados de una cuenca de igual profundidad y dimensiones en otra ciudad, sabe que un pulso transversal demorará $2.50\text{ s}$ en cruzar la cuenca planificada por el arquitecto. Demuestre al arquitecto que habrá varias resonancias de seiche posibles en la cuenca hidrológica para las bajas frecuencias típicas de los terremotos en el rango de $0 - 4\text{ Hz}$.